\section{Introduction}
\label{intro}

Reward hacking is a fundamental challenge in reinforcement learning (RL) where
an agent exploits loopholes in the reward function to achieve high rewards
without performing the intended task~\cite{rewardhacking}. This phenomenon can
lead to unintended consequences and suboptimal performance, posing a significant
threat to the safet and reliability of AI systems~\cite{anthropicrh}. As the
state of the art AI models continue to be post-trained in RL enovironements, the
risk of reward hacking becomes increasingly pertinent. 


In this work, we examine the propensity of reward hacking in programming tasks. We observe that when provided test cases, models will reward hack (eg. hardcode the test cases) 20\% of the time. We propose a 
potential solution: formal specifications~\cite{Verus}. Formal specfications provide a mathematical guarantee that a program obeys a specification for all available inputs. If a program satisfies a 
correct formal specification, it is guaranteed to be correct for all inputs, and thus reward hacking is not possible. While in practice,
formal specifications are difficult to write and underspecifications are
theoretically reward hackable. However, we empirically observe that formal 
specifications were rewawrd hacked 0\% of the time.

While this paints a rosy picture, there are still major challenges with formal verification. Formal languages are more difficult for LLMs to generate and places an additional burden: they must write a proof that their code satisfies the specification. This shows up in our results: LLMs are able to generate correct implementations for 70\% of the benchmarks when given test cases, but only 21\% when given formal specifications. Thus, while formal specifications may be a promising solution to reward hacking, there are still significant challenges to overcome before they can be widely adopted.



Our contributions are as follows:
\begin{itemize}
\item We propose formal specifications as a potential solution to reward hacking in programming tasks.
\item We develop a measure of reward hacking in programming tasks and use it to empirically evaluate the prevalence of reward hacking when using test cases versus formal specifications.
\item We find that reward hacking occurs 20\% of the time with test cases and 0\% of the time with formal specifications and discuss future work to address the challenges of using formal specifications in practice.
\end{itemize}

\subsection{Background}
\label{background}

A \emph{formal specification} mathematically describes the expected properties of a program. 
Unlike testing, which checks a finite set of inputs, 
formal specifications can quantify over all possible inputs, 
enabling them to provide guarantees over all possible program behaviors.
\emph{Formal verification} is the process of mathematically proving that a program satisfies its specification.
Thus, it shifts the burden from ensuring correctness of an entire implementation 
(which may have complex performance optimizations)
to ensuring correctness of a higher-level mathematical specification, 
an easier task. 

% typically using a verification tool to ensure proof correctness.

One state-of-the-art verification tool is Verus~\cite{Verus}, a semi-automated Rust verifier which 
has already been used at the intersection
of AI and verification \cite{AutoVerus,alphaverus,verussage,RAGVerus,vericode}. 
In Verus, users can write three kinds of code: 
\begin{itemize}
    \item \emph{spec} code describes higher-level program properties such as function pre- and post-conditions
    or loop invariants.
    \item \emph{proof} annotations help give Verus hints for proving an implementation satisfies its spec:
    Verus is semi-automated, meaning it can automatically prove many specifications, 
    but when more complex reasoning is required,
    it may need some help in knowing intermediate proof goals 
    or useful lemma calls.
    \item \emph{exec} code is simply regular Rust executable code. 
    During compilation, Verus erases spec and proof code, 
    leaving only the exec code: 
    there is no runtime performance cost associated with adding Verus proofs.
\end{itemize}
Verus is designed to integrate directly into Rust code: 
it relies on the Rust type-checker and uses Rust macros and annotations.
However, using Verus does involve some Verus-specific syntax, 
such as additional mathematical keywords or naming function outputs to describe their expected properties.

% Verus will then check if the code satisfies these requirements. 
% Note that Verus's verification result is only as good as the specification quality: 
% if the specification is incorrect or fails to capture all possible cases, 
% implementations verified against such specifications will also lack full guarantees.
% And since Verus specifications are checkable and compositional, they can serve as unambiguous descriptions of program intent---beyond what unit tests can capture. 
% Verus~\cite{Verus} has already been used in projects at the intersection
% of AI and verification \cite{AutoVerus,alphaverus,verussage,RAGVerus,vericode}. 